Against the first coding, the model marks fewer elements, and most element shares are lower. In a later check against the coder's re-code of 40 paragraphs, made before the model run, the recall deficit persists in point estimate but is smaller (Section~\ref{subsec:directional}). Opposite errors can partly cancel in aggregate figures, whereas one-directional errors shift them \parencite[p.~2091]{edwards2026rogangladen}. Aggregate shares therefore need a defended human reference. Sections~\ref{subsec:interpretation} and~\ref{subsec:mechanisms} interpret the pattern and weigh possible sources, none isolated by the design. Sections~\ref{subsec:theoretical_implications} to~\ref{subsec:limitations_future} cover implications, limitations and next steps.

\subsection{The Recall Deficit Depends on the Reference Reading}\label{subsec:interpretation}\label{subsec:comparison_prior}

In a later split against the first coding, element precision exceeds recall (Figure~\ref{fig:overview}b), whereas recall exceeds precision in 20 of 27 codebook-prompted task replications by \textcite{pangakis2025humansintheloop}. The higher recall against the re-code made before the model run shows this asymmetry to be reference-sensitive, not an established measure of model-side omission.

The gap recurs in most configurations, so the shared setup remains a candidate explanation. Yet a later fresh-census check finds the gap small among valid outputs, a selected subset (Section~\ref{subsec:recurrence}). Against the evaluation sample's first coding, a later Opus~5 run is near balance in point estimate, so the setup does not force a deficit.

\subsection{Possible Sources of Element Disagreement}\label{subsec:mechanisms}

Against the first coding, most first-execution misses occur on technology rows where the model marked another element (later check, Tables~\ref{tab:fn_causes} and~\ref{tab:checks_4p}). The union of all three registered executions raises recall only from 0.557 to 0.596, so the disagreements persist across executions. Persistence does not establish which side is wrong.

On the reference side, the second coding holds fewer element tuples than the first (Table~\ref{tab:added_analyses}). This may reflect threshold drift, a known risk of coding over time \parencite[p.~565]{song2020validations}, learning, memory or, after the results, knowledge of them, which the design cannot separate.

Two mechanisms from the literature, both conjectures here, could produce a deficit of this sign. A model asked for several labels may stop listing early, as autoregressive models tended to concentrate probability on one label per step \parencite[p.~2472]{ma2025multilabel}. Label bias, a preference for some answers regardless of the input, may also favour absence \parencite{reif2024labelbias}. Both studies examine open-weight models' output probabilities on other tasks, but Sonnet~4.6 returns only text. The mandatory form's recall gain, clear only on the fresh census and confounded with a doubled output cap, fits the first conjecture without isolating it (Table~\ref{tab:form_variant}). Under the first conjecture, Mercury~3, a diffusion model that does not generate strictly left to right, should show a narrower gap than autoregressive models, as declared before its runs \parencite{inception2026mercury3,inception2025mercury}. Its gap was larger than the registered model's, though model, temperature and provider also differed (Table~\ref{tab:pr_gaps}). A third candidate, ambiguous definitions, is not singled out: in the perturbation battery, the model's labels move with definition wording but also with frame order, element labels and a negative-control sentence (Appendix~\ref{app:battery}).

The coder's non-blind re-reading of misses points to both model and reference (Section~\ref{subsec:missed_results}). Its judgements are not error rates. Of the 13 sampled investment misses, 5 were judged possibly over-marked. Four of these report deployment growth or decline, a separate element (registered re-code $\alpha$ 0.369) that the model marked each time. Investment's comparatively high stability on shared rows thus does not show that its coding follows the intended boundary (Table~\ref{tab:alpha_waves}).

\subsection{Theoretical Implications}\label{subsec:theoretical_implications}\label{subsec:implications}

The instrument measures presence, not salience, and most elements can fire on a topical reference alone. The residual frame appears on two evaluation paragraphs in the first coding and on six of 55 in the census reference (Tables~\ref{tab:displacement} and~\ref{tab:fresh_displacement}). Neither count shows that the inventory is complete.

Measuring frames through elements \parencite{ali2025hydrogenframing,liu2023strategyframing} exposed differences a single label would hide. Within the feasibility frame, F1 is 0.723 for technical feasibility and 0.433 for technical performance, yet their later intervals overlap (Table~\ref{tab:checks_4o}). Technical performance is unstable in both re-code waves, so the contrast is tentative (Table~\ref{tab:alpha_waves}).

Labels that discriminate well can still yield biased shares when counted \parencite[pp.~7--8]{linder2026multicalibrated}, so label and prevalence agreement remain distinct. The later design-based adjustment moves shares towards the convention of its human-coded sample, whose accuracy its authors name as a condition \parencite[Sec.~6]{egami2023dsl}, and does not correct reference error (Section~\ref{subsec:dsl_results}). Under the second coding, made mostly after the results, it gives an exploratory sensitivity, not a better estimate. Label agreement, reference stability and derived shares thus need joint validation.

\subsection{Managerial Implications}\label{subsec:managerial_implications}

Media-monitoring teams can report the technologies detected in prefiltered paragraphs, and the evaluations of frames a coder confirms, once both are validated per language on their own material. Whole technology records rarely match exactly, so exporting them without review is not supported (Section~\ref{subsubsec:categorical}).

Element, frame and evaluation shares should not be reported without a human-coded probability sample. A blind re-code completed before model results are inspected assesses the coder's temporal stability. Re-reading sampled disagreements can indicate whether model or reference may err. Once the reference is secured, shares need adjustment towards a defended convention where warranted. A second coder can help defend the convention, but no step establishes validity by itself. This workflow transfers to other monitoring projects, which must define their own frames and elements and revalidate.

Each source, language, country, technology or period to be compared needs its own human-coded sample. When analysts compare frame shares across outlets, countries and periods, misses that vary by group can exaggerate, hide or reverse the difference human coding would show \parencite[p.~2092]{edwards2026rogangladen}. Against the first coding, German technology detection was lower end to end, mostly on paragraphs whose first execution failed. Among valid outputs, the contrast is inconclusive (Section~\ref{subsec:agreement}). All five such failures were German, but the fresh census had four per language (Table~\ref{tab:checks_4o}). With element displacement intervals 5.6 to 17.0 points wide, the present sample is too small per cell (Table~\ref{tab:topic_requirements}).

Model calls scale to larger corpora, but sourcing and validation do not yet. The registered run's 360 executions took 12.5 minutes and cost an estimated USD 2.18 in model calls. Coding and validation labour was not measured. Larger applications need a licensed or API-based source (Table~\ref{tab:topic_requirements}). They should budget retries, report attrition per language, fix configurations, journal every call and revalidate new model versions (Section~\ref{subsec:recurrence}).

\subsection{Limitations and Future Research}\label{subsec:limitations_future}\label{subsec:limitations}

Both human readings come from one coder, the author, and neither is independently validated. Calibration covered 32 paragraphs on one day. The final split of deployment growth from investment was not trial-coded before the freeze. The registered re-code came after as few as five days instead of 14, so memory may inflate its temporal stability and sustain the deficit against it. Imperfect human labels can also depress estimated classification performance, as simulations show \parencite[p.~560]{song2020validations}.

The supplementary re-code followed the registered results and the re-reading of misses, so it may have moved towards the model. Omitting the 25 re-read paragraphs does not remove knowledge of the results and leaves a non-random subset (Appendix~\ref{app:later_checks}). Model-agreement and prevalence intervals condition on the chosen labels and do not propagate uncertainty about alternative human readings. A discarded first attempt with keyword suggestions may have shaped technology ticks (Section~\ref{subsec:annotation}). The general category's figures also depend on its reading.

Model and coder differed in instructions (Section~\ref{subsec:codebook}) and in context. The coder saw headline, outlet, language and often other article paragraphs. The model saw only the paragraph, although the codebook's unit definition names headline and metadata. Context could reveal elements despite the shared paragraph-only rule and so inflate the apparent deficit. Adding headline, outlet and language in one later execution per sample did not consistently raise recall or narrow the gap (Table~\ref{tab:context_variant}) but left the coder's knowledge of other paragraphs untested. With one coder, mostly single runs of further models and variation of only the model's context, the deficit cannot be apportioned among model, codebook, reference and context.

The reference and context limits point to the next test. Two independent coders would apply a narrower, untested codebook to the model's paragraph-only input on fresh paragraphs. An element would require a technology-anchored claim (not a topic mention), a stated boundary to its neighbours and a rule on whether one cue may support two elements. The general category would need one explicit reading. Procedural reporting would need a settled boundary to regulation framing. Candidate rules are in Table~\ref{tab:battery_sharpen}. A model-assisted review of drafts may find gaps early (Table~\ref{tab:development_review}).

The sample of 120 paragraphs from 56 articles is small and clustered: the element micro-F1 interval spans about $\pm 0.06$ and per-language intervals up to $\pm 0.09$. Population estimates target the eligible frame of a purposive, hand-collected corpus, defined by a keyword prefilter with unmeasured misses. The later analyses (Table~\ref{tab:added_analyses}) and appendix checks mostly reuse inspected samples and report intervals unadjusted for multiplicity, so none is confirmatory. The evidence covers one topic, one codebook and two languages from different outlets, not other frames or domains.

The data rule restricts article text to the registered run's EU route. Exceptions were authorised for several later steps. Mercury~3 received evaluation, census, calibration and development paragraphs via its developer's interface, hosted at an unknown location. After the EU quota ran out, 62\% of battery executions used a global route, possibly outside the EU. Assistant agents outside the EU route read paragraphs to derive the battery's candidate rules. A language-model assistant consulted during the fresh census received census text but recommended no label, and none was changed (Appendix~\ref{app:reference_disclosures}). Some assistant data routes cannot be reconstructed from the archived material (Appendix~\ref{app:guide_llm}). The records document no rights basis for storing or externally processing the hand-collected texts, so they are withheld as a precaution (Table~\ref{tab:topic_requirements}).

Proprietary Sonnet~4.6 exposes no seed. Its provider may change or end access to this version \parencite{feuerriegel2026guidellm}. The coding interface's build is unarchived.

As the model comparisons reused inspected samples, production selection needs a fresh selection sample, a rule fixed beforehand and a separate, untouched evaluation sample.

The topic description's broader project asks how energy coverage changes around the 2024 United States presidential election. It needs a human-coded probability sample per country, language and period, a sample of keyword-rejected paragraphs to measure the filter's misses, and source, model version and data route fixed for the study period.
